% MDPI generic article skeleton.
% For a real submission, download the journal-specific class file from the target MDPI journal's "Instructions for Authors" page (e.g., Definitions/mdpi.cls), drop it next to main.tex, and switch the documentclass line below.

\documentclass[a4paper,10pt]{article}
\usepackage[a4paper,margin=2cm]{geometry}

\usepackage{amsmath,amssymb,amsfonts}
\usepackage{graphicx}
\usepackage{booktabs}
\usepackage{tabularx}
\usepackage{cite}
\usepackage{xcolor}
\usepackage{hyperref}
\usepackage{url}

\title{Paper Title}
\author{First Author$^{1}$, Second Author$^{2,*}$\\
$^{1}$Department, Institution, City, Country\\
$^{2}$Department, Institution, City, Country\\
$^{*}$Corresponding author: \texttt{email@example.com}}
\date{}

\begin{document}

\maketitle

\begin{abstract}
The abstract goes here. MDPI journals usually allow up to 200 words for research articles and up to 300 for reviews. Summarise motivation, methods, results, and significance.
\end{abstract}

\noindent\textbf{Keywords:} keyword one, keyword two, keyword three

\section{Introduction}
\label{sec:introduction}

Introduction text. State the problem, the gap in prior work, and the contribution of this paper.

\section{Related Work}
\label{sec:related}

\section{Materials and Methods}
\label{sec:method}

\section{Results}
\label{sec:results}

\section{Discussion}
\label{sec:discussion}

\section{Conclusions}
\label{sec:conclusion}

\bibliographystyle{plain}
\bibliography{references}

\end{document}
